\ifdefined\gkver\else\def\gkver{student}\fi
\documentclass[\gkver]{gaokaozhenti}
\begin{document}

\chapter{2013年四川理综}
%% paperType: 理综物理部分

\begin{enumerate}
\item
%% number: 1
%% paperName: 2013年四川理综第 1 题 6 分
%% typeId: 1
%% type: 单选题
%% chapter: 光学
%% point: 光的电磁说
%% method: 无
%% score: 6
%% degree: 850
%% duplicateId: 0
%% body:
下列关于电磁波说法，正确的是\xzanswer{C}

\fourchoices[answer=C]
{电磁波只能在真空中传播}
{电场随时间变化时一定产生电磁波}
{做变速运动的电荷会在空间产生电磁波}
{麦克斯韦第一次用实验证实了电磁波的存在}

%% answer: C
%% memo:
\memoanswer{
本题原卷及材料中未提供详解，待补充。
}

\item
%% number: 2
%% paperName: 2013年四川理综第 2 题 6 分
%% typeId: 1
%% type: 单选题
%% chapter: 交流电
%% point: 交流电
%% method: 无
%% score: 6
%% degree: 850
%% duplicateId: 0
%% body:
用 $220\UV$ 的正弦交流电通过理想变压器对一负载供电，变压器输出电压是 $110\UV$，通过负载的电流图象如图所示，则\xzanswer{A}

\onepicture[w=7cm, align=right]{figs/02.png}

\fourchoices[answer=A]
{变压器输入功率约为 $3.9\UW$}
{输出电压的最大值是 $110\UV$}
{变压器原、副线圈的匝数比是 $1:2$}
{负载电流的函数表达式 $i=0.05\sin\left(100\pi t+\frac{\pi}{2}\right)\UA$}

%% answer: A
%% memo:
\memoanswer{
本题原卷及材料中未提供详解，待补充。
}

\item
%% number: 3
%% paperName: 2013年四川理综第 3 题 6 分
%% typeId: 1
%% type: 单选题
%% chapter: 光学
%% point: 光的折射
%% method: 无
%% score: 6
%% degree: 650
%% duplicateId: 0
%% body:
光射到两种不同介质的分界面，分析其后的传播形式可知\xzanswer{D}

\fourchoices[answer=D]
{折射现象的出现表明光是纵波}
{光总会分为反射光和折射光}
{折射光与入射光的传播方向总是不同的}
{发生折射是因为光在不同介质中的传播速度不同}

%% answer: D
%% memo:
\memoanswer{
本题原卷及材料中未提供详解，待补充。
}

\item
%% number: 4
%% paperName: 2013年四川理综第 4 题 6 分
%% typeId: 1
%% type: 单选题
%% chapter: 曲线运动
%% point: 天体运动
%% method: 无
%% score: 6
%% degree: 650
%% duplicateId: 0
%% body:
太阳系外行星大多不适宜人类居住，绕恒星“Glicsc581”运行的行星“Gl-581c”却很值得我们期待。该行星的温度在 $0\UdC$ 到 $40\UdC$ 之间，质量是地球的 6 倍，直径是地球的 1.5 倍、公转周期为 13 个地球日。“Glicsc581”的质量是太阳质量的 0.31 倍。设该行星与地球均视为质量分布均匀的球体，绕其中心天体做匀速圆周运动，则\xzanswer{B}

\fourchoices[answer=B]
{在该行星和地球上发射卫星的第一宇宙速度相同}
{如果人到了该行星，其体重是地球上的 $2\frac{2}{3}$ 倍}
{该行星与“Glicsc581”的距离是日地距离的 $\sqrt{\frac{365}{13}}$ 倍}
{由于该行星公转速率比地球大，地球上的米尺如果被带上该行星，其长度一定会变短}

%% answer: B
%% memo:
\memoanswer{
本题原卷及材料中未提供详解，待补充。
}

\item
%% number: 5
%% paperName: 2013年四川理综第 5 题 6 分
%% typeId: 1
%% type: 单选题
%% chapter: 机械振动
%% point: 振动图象
%% method: 无
%% score: 6
%% degree: 650
%% duplicateId: 0
%% body:
图 1 是一列简谐横波在 $t=1.25\Us$ 时的波形图，已知 c 位置的质点比 a 位置的晚 $0.5\Us$ 起振。则图 2 所示振动图像对应的质点可能位于\xzanswer{D}

\twopicture[numstyle=arabic, numA=1, numB=2, w=6cm, align=right]
{figs/05a.png}
{figs/05b.png}

\fourchoices[answer=D]
{$a<x<b$}
{$b<x<c$}
{$c<x<d$}
{$d<x<e$}

%% answer: D
%% memo:
\memoanswer{
本题原卷及材料中未提供详解，待补充。
}

\item
%% number: 6
%% paperName: 2013年四川理综第 6 题 6 分
%% typeId: 2
%% type: 多选题
%% chapter: 直线运动
%% point: 运动图象
%% method: 无
%% score: 6
%% degree: 650
%% duplicateId: 0
%% body:
甲、乙两物体在 $t=0$ 时刻经过同一位置沿 $x$ 轴运动，其 $v-t$ 图像如图所示。则\xzanswer{BD}

\onepicture[w=8cm, align=right]{figs/06.png}

\fourchoices[answer=BD]
{甲、乙在 $t=0\Us$ 到 $t=1\Us$ 之间沿同一方向运动}
{乙在 $t=0$ 到 $t=7\Us$ 之间的位移为零}
{甲在 $t=0$ 到 $t=4\Us$ 之间做往复运动}
{甲、乙在 $t=6\Us$ 时的加速度方向相同}

%% answer: BD
%% memo:
\memoanswer{
根据 $v-t$ 图象可知，在 $t=0$ 到 $t=1\Us$ 这段时间内，甲一直向正方向运动，而乙先向负方向运动再向正方向运动，故 A 选项错误；根据 $v-t$ 图象可知，在 $t=0$ 到 $t=7\Us$ 这段时间内，乙的位移 $x=-\frac{v_{0}}{2}\times0.5\Us+\frac{v_{0}}{2}\times0.5\Us+\frac{v_{0}}{2}\times3\Us-\frac{v_{0}}{2}\times3\Us=0$，故 B 选项正确；根据 $v-t$ 图象可知，在 $t=0$ 到 $t=4\Us$ 这段时间内，甲一直向正方向运动，故 C 选项错误；根据 $v-t$ 图象可知，在 $t=6\Us$ 时，甲、乙的加速度方向均与正方向相反，故 D 选项正确。
}

\item
%% number: 7
%% paperName: 2013年四川理综第 7 题 6 分
%% typeId: 2
%% type: 多选题
%% chapter: 电磁感应
%% point: 电磁感应电路分析
%% method: 无
%% score: 6
%% degree: 450
%% duplicateId: 0
%% body:
如图所示，边长为 $L$、不可形变的正方形导体框内有半径为 $r$ 的圆形区域，其磁感应强度 $B$ 随时间 $t$ 的变化关系为 $B=kt$（常量 $k>0$）。回路中滑动变阻器 $R$ 的最大阻值为 $R_{0}$，滑动片 P 位于滑动变阻器中央，定值电阻 $R_{1}=R_{0}$、$R_{2}=\frac{R_{0}}{2}$。闭合开关 S，电压表的示数为 $U$，不考虑虚线 MN 右侧导体的感应电动势。则\xzanswer{AC}

\onepicture[w=6cm, align=right]{figs/07.png}

\fourchoices[answer=AC]
{$R_{1}$ 两端的电压为 $\frac{U}{7}$}
{电容器的 a 极板带正电}
{滑动变阻器 R 的热功率为电阻 $R_{2}$ 的 5 倍}
{正方形导线框中的感应电动势为 $kL^{2}$}

%% answer: AC
%% memo:
\memoanswer{
本题原卷及材料中未提供详解，待补充。
}

\item
%% number: 8
%% paperName: 2013年四川理综第 8 题 11 分
%% typeId: 4
%% type: 实验题
%% chapter: 功和能
%% point: 动能定理
%% method: 无
%% score: 11
%% degree: 450
%% duplicateId: 0
%% body:
\begin{enumerate}
	\item
	在探究两电荷间相互作用力的大小与哪些因素有关的实验中，一同学猜想可能与两电荷的间离和带电量有关。他选用带正电的小球 A 和 B，A 球放在可移动的绝缘座上，B 球用绝缘丝线悬挂于玻璃棒 C 点，如图所示。

	实验时，先保持两电荷量不变，使 A 球从远处逐渐向 B 球靠近，观察到两球距离越小，B 球悬线的偏角越大，再保持两球的距离不变，改变小球所带的电荷量，观察到电荷量越大，B 球悬线的偏角越大。

	实验表明：两电荷之间的相互作用力，随其距离的\tkanswer[1.2cm]{减小}而增大，随其所带电荷量\tkanswer[1.2cm]{增大}而增大。

	此同学在探究中应用的科学方法是\tkanswer[2.6cm]{控制变量法}（选填：“累积法”、“等效替代法”、“控制变量法”、或“演绎法“）。

	\onepicture[w=5cm, align=right]{figs/08.png}

	\item
	如图 1 所示，某组同学借用“探究 $a$ 与 $F$、$m$ 之间的定量关系”的相关实验思想、原理及操作，进行“探究合外力做功和动能变化的关系”的实验：

	\begin{steps}
		\item
		为达到平衡阻力的目的，取下细绳和托盘，通过调节垫片的位置，改变长木板倾斜程度，根据打出的纸带判断小车是否做\tkanswer[2cm]{匀速直线}运动。

		\item
		连接细绳及托盘，放入砝码，通过实验得到图 2 所示的纸带。纸带上 O 为小车运动起始时刻所打的点，选取时间间隔为 $0.1\Us$ 的相邻计数点 A、B、C、D、E、F、G。实验时小车所受拉力为 $0.2\UN$，小力车的质量为 $0.2\Ukg$。

		请计算小车所受合外力做的功 $W$ 和小车动能的变化 $\Delta E_{\mathrm{k}}$，补填表中空格（结果保留至少小数点后第四位）。

		\begin{center}
		\begin{tabular}{|c|c|c|c|c|c|}
		\hline
		 & O-B & O-C & O-D & O-E & O-F \\ \hline
		$W/\UJ$ & 0.0432 & 0.0572 & 0.0734 & 0.0915 & \\ \hline
		$\Delta E_{\mathrm{k}}/\UJ$ & 0.0430 & 0.0570 & 0.0734 & 0.0907 & \\ \hline
		\end{tabular}
		\end{center}

		分析上述数据可知：在实验误差允许的范围内 $W=\Delta E_{\mathrm{k}}$，与理论推导结果一致。

		\item
		实验前已测得托盘的质量 $7.7\times10^{-3}\Ukg$，实验时该组同学放入托盘中的砝码质量应为\tkanswer[1.5cm]{0.015}$\Ukg$（$g$ 取 $9.8\Umsq$，结果保留至小数点后第三位）。

	\end{steps}

	\onepicture[w=8cm, align=right]{figs/08a.png}

	\onepicture[w=11cm, align=right]{figs/08b.png}

\end{enumerate}

%% answer: 
\jdanswer{
\begin{enumerate}
	\item
	减小，增大，控制变量法

	\item
	匀速直线运动（3 分）；0.1115（2 分），0.1105（3 分）；0.015（3 分）

\end{enumerate}
}
%% memo:
\memoanswer{
本题原卷及材料中未提供详解，待补充。
}

\item
%% number: 9
%% paperName: 2013年四川理综第 9 题 15 分
%% typeId: 6
%% type: 计算题
%% chapter: 运动和力
%% point: 牛顿第二定律
%% method: 无
%% score: 15
%% degree: 650
%% duplicateId: 0
%% body:
近来，我国多个城市开始重点治理“中国式过马路”行为。每年全国由于行人不遵守交通规则而引发的交通事故上万起，死亡上千人。只有科学设置交通管制，人人遵守交通规则，才有保证行人的生命安全。

如右图所示，停车线 AB 与前方斑马线边界 CD 间的距离为 $23\Um$。质量 $8\Ut$、车长 $7\Um$ 的卡车以 $54\Ukmh$ 的速度向北匀速行驶，当车前端刚驶过停车线 AB，该车前方的机动车交通信号灯由绿灯变黄灯。
\begin{enumerate}
	\item
	若此时前方 C 处人行横道路边等待的行人就抢先过马路，卡车司机发现行人，立即制动，卡车受到的阻力为 $3\times10^{4}\UN$。求卡车的制动距离？
	\item
	若人人遵守交通规则，该车将不受影响地驶过前方斑马线边界 CD。为确保行人安全，D 处人行横道信号灯应该在南北向机动车信号灯变黄灯后至少多久变为绿灯？
\end{enumerate}

\twopicture[showlabel=false, h=4cm, align=right]
{figs/09a.png}
{figs/09b.png}

%% answer: 
\jdanswer{
\begin{enumerate}
	\item
	$30\Um$

	\item
	$2\Us$

\end{enumerate}
}
%% memo:
\memoanswer{
卡车质量 $m=8\times10^{3}\Ukg$，初速度 $v_{0}=15\Ums$。

（1）设卡车的制动距离为 $s_{1}$，由动能定理得：$-fs_{1}=0-\frac{1}{2}mv^{2}$，$s_{1}=30\Um$。

（2）$s_{0}+l=v_{0}\Delta t$，

即：$23+7=15\Delta t$，$\therefore\Delta t=2\Us$。
}

\item
%% number: 10
%% paperName: 2013年四川理综第 10 题 17 分
%% typeId: 6
%% type: 计算题
%% chapter: 运动和力
%% point: 牛顿定律的应用
%% method: 无
%% score: 17
%% degree: 450
%% duplicateId: 0
%% body:
在如图所示的竖直平面内，物体 A 和带正电的物体 B 用跨过定滑轮的绝缘轻绳连接，分别静止于倾角 $\theta=37^{\circ}$ 的光滑斜面上的 M 点和粗糙绝缘水平面上，轻绳与对应平面平行。劲度系数 $k=5\UNm$ 的轻弹簧一端固定在 O 点，一端用另一轻绳穿过固定的光滑小环 D 与 A 相连，弹簧处于原长，轻绳恰好拉直，DM 垂直于斜面。水平面处于场强 $E=5\times10^{4}\UNC$、方向水平向右的匀强电场中。已知 A、B 的质量分别为 $m_{A}=0.1\Ukg$，$m_{B}=0.2\Ukg$，B 所带电荷量 $q=+4\times10^{-6}\UC$。设两物体均视为质点，不计滑轮质量和摩擦，绳不可伸长，弹簧始终在弹性限度内，B 电量不变。取 $g=10\Umsq$，$\sin37^{\circ}=0.6$，$\cos37^{\circ}=0.8$。
\begin{enumerate}
	\item
	求 B 所受摩擦力的大小；
	\item
	现对 A 施加沿斜面向下的拉力 $F$ 使 A 以加速度 $a=0.6\Umsq$ 开始作匀加速直线运动。A 从 M 到 N 的过程中，B 的电势能增加了 $\Delta E_{p}=0.06\UJ$。已知 DN 沿竖直方向，B 与水平面间的动摩擦因数为 $\mu=0.4$。求 A 到达 N 点时拉力 $F$ 的瞬时功率？
\end{enumerate}

\onepicture[w=6cm, align=right]{figs/10.png}

%% answer: 
\jdanswer{
\begin{enumerate}
	\item
	$0.4\UN$

	\item
	$0.528\UW$

\end{enumerate}
}
%% memo:
\memoanswer{
（1）对 A：绳拉力 $T=m_{A}g\sin\theta=0.6\UN$；

对 B：受电场力 $F=qE=0.2\UN<T$；

故 B 受到的摩擦力：$f=T-F=0.4\UN$。

（2）A 从 M 到 N 的过程中，B 的电势能增加了 $\Delta E_{p}=0.06\UJ$，则克服电场力做功为 $W=qEs$，$\therefore s=0.3\Um$。

A、B 匀加速直线运动，加速度大小 $a$、运动的位移大小 $s$、末速度大小 $v$ 都相等，$v=\sqrt{2as}=0.6\Ums$。

对 B：$T_{1}-qE-\mu m_{B}g=m_{B}a$，$T_{1}=1.12\UN$。

对 A 到达 N 点时：弹簧的伸长量 $x=ON-OM=\frac{s}{\sin\theta}-s\cot\theta=0.1\Um$。

弹力 $F_{2}=kx=5\times0.1\UN=0.5\UN$。

由牛顿第二定律得：$F+m_{A}g\sin\theta-T_{1}-F_{2}\sin\theta=m_{A}a$，$F=0.88\UN$。

A 到达 N 点时拉力 $F$ 的瞬时功率 $P=Fv=0.88\times0.6\UW=0.528\UW$。
}

\item
%% number: 11
%% paperName: 2013年四川理综第 11 题 19 分
%% typeId: 6
%% type: 计算题
%% chapter: 磁场
%% point: 带电粒子在磁场中的运动
%% method: 无
%% score: 19
%% degree: 450
%% duplicateId: 0
%% body:
如图所示，竖直平面（纸面）内有平面直角坐标系 xOy，x 轴沿水平方向。在 $x\leq0$ 的区域内存在方向垂直纸面向里，磁感应强度大小为 $B_{1}$ 的匀强磁场。在第二象限紧贴 y 轴固定放置长为 $l$、表面粗糙的不带电绝缘平板，平板平行 x 轴且与 x 轴相距 $h$。在第一象限内的某区域存在方向互相垂直的匀强磁场（磁感应强度大小为 $B_{2}$，方向垂直于纸面向外）和匀强电场（图中未画出）。一质量为 $m$、不带电的小球 Q 从平板下侧 A 点沿 x 正向抛出；另一质量也为 $m$、带电量为 $q$ 的小球 P 从 A 点紧贴平板沿 x 轴正向运动，变为匀速运动后从 y 轴上的 D 点进入电磁场区域做匀速圆周运动，经 $\frac{1}{4}$ 圆周离开电磁场区域，沿 y 轴负方向运动，然后从 x 轴上的 K 点进入第四象限。小球 P、Q 相遇在第四象限内的某一点，且竖直方向速度相同。设运动过程中小球 P 的电量不变，小球 P 和 Q 始终在纸面内运动且均看作质点，重力加速度为 $g$。求：
\begin{enumerate}
	\item
	匀强电场的场强大小，并判断 P 球所带电荷的正负；
	\item
	小球 Q 的抛出速度 $v_{0}$ 取值范围；
	\item
	$B_{1}$ 是 $B_{2}$ 的多少倍？
\end{enumerate}

\onepicture[w=7cm, align=right]{figs/11.png}

%% answer: 
\jdanswer{
\begin{enumerate}
	\item
	$\frac{mg}{q}$，正电

	\item
	$0<v_{0}\leq\left(\frac{m^{2}g}{q^{2}B_{1}B_{2}}+l\right)\sqrt{\frac{g}{2h}}$

	\item
	0.5

\end{enumerate}
}
%% memo:
\memoanswer{
（1）由题意，小球 P 到达 D 前匀速：$qvB_{1}=mg$，$v=\frac{mg}{qB_{1}}$ ①，且 P 带正电。

过 D 进入电磁场区域做匀速圆周运动，则 $mg=qE$，

$E=\frac{mg}{q}$ ②。

（2）经 $\frac{1}{4}$ 圆周离开电磁场区域，沿 y 轴负方向运动，设匀速圆周运动半径为 $R$，

则 $qvB_{2}=m\frac{v^{2}}{R}$ ③。

设小球 P、Q 在第四象限相遇时的坐标为 $x$、$y$。

设小球 Q 平抛后经时间 $t$ 与 P 相遇，Q 的水平位移为 $s$，竖直位移为 $d$。则

$s=v_{0}t$ ④，$d=\frac{1}{2}gt^{2}$ ⑤。

位移关系：$x=s-l=R$ ⑥，$y=-(d-h)\leq0$ ⑦。

\onepicture[w=8cm, align=right]{figs/11_an.png}

联解①③④⑤⑥⑦得：$0<v_{0}\leq\left(\frac{m^{2}g}{q^{2}B_{1}B_{2}}+l\right)\sqrt{\frac{g}{2h}}$ ⑧。

（3）小球 Q 在空间作平抛运动，在满足题目条件下，运动到小球 P 穿出电磁场时的高度（图中 N 点）时，设用时 $t_{1}$，竖直方向速度 $v_{y}$，竖直位移 $y_{Q}$，则有

$y_{Q}=R=\frac{1}{2}gt_{1}^{2}$ ⑨，

$v_{y}=v=gt_{1}$ ⑩ （此后 $a=g$ 同，PQ 相遇点竖直方向速度必同）。

联解得：$B_{1}=0.5B_{2}$。
}

\end{enumerate}

\end{document}
